\subsection{Emergencia conjunta de las cinco construcciones}
\label{sec:cinco-construcciones-continuo-ch}

Para \(0\leq k\leq N\), sea
\[
 \rho_{N,k}:=\rho_k\circ\rho_{k+1}\circ\cdots\circ\rho_{N-1},
 \qquad \rho_{k,k}:=\operatorname{id}_{H_k},
\]
y, para una historia superviviente \(\zeta\in H_k\), defínase
\[
 \operatorname{Desc}_{k,N}(\zeta)
 :=\{x\in H_N:\rho_{N,k}(x)=\zeta\}.
\]
El árbol finito \(\mathscr T_{k,N}(\zeta)\) tiene como vértices
\[
 \coprod_{j=k}^{N}\{y\in H_j:\rho_{j,k}(y)=\zeta\}
\]
y como aristas únicamente las prolongaciones TPK
\(a:y\to z\) con \(z\in\rho_j^{-1}(y)\). Por tanto, ni sus vértices ni
sus prolongaciones proceden de un árbol auxiliar.

Para \(x\in\operatorname{Desc}_{k,N}(\zeta)\), la rama marcada
\(\lambda_{\zeta,x}\) es la sucesión
\[
 \zeta=x_k\longrightarrow x_{k+1}\longrightarrow\cdots
 \longrightarrow x_N=x
\]
determinada por los truncamientos de \(x\). El registro base común es
\[
\begin{aligned}
 \mathscr R_{k,N}(\zeta;x):=
 \bigl(&\mathscr T_{k,N}(\zeta),\lambda_{\zeta,x};\\
 &\bigl(\varepsilon_a,s_a,o_a,
   (\delta_{\diamond,a},r_{\diamond,a},q_{\diamond,a})
       _{\diamond\in\{\Sigma,\Pi\}},
   \operatorname{Carr}_a,C_a,t_a,\operatorname{Fr}_a,
   \operatorname{Ruta}_a,\operatorname{Mem}_a,
   \operatorname{Surv}_a\bigr)_a;\\
 &\operatorname{src},\operatorname{tgt},\circ,
   (\rho_j)_{k\leq j<N}\bigr),
\end{aligned}
\tag{C1}
\label{eq:base-comun-registros}
\]
donde
\[
 U_a(m_0)=C_am_0+t_a,
 \qquad
 \delta_{\diamond,a}=r_{\diamond,a}+9q_{\diamond,a}.
\]
El símbolo \(m_0\) de esta fórmula es una coordenada de memoria; no es la
profundidad modular \(m\) utilizada más abajo. El espacio base verdadero es
\[
 \mathfrak R_{k,N}(\zeta)
 :=\{\mathscr R_{k,N}(\zeta;x):
       x\in\operatorname{Desc}_{k,N}(\zeta)\}.
\]
La rama marcada no es una entrada añadida: es la sucesión de aristas
efectivamente producidas que termina en \(x\). El resto del registro conserva
simultáneamente las fibras laterales del árbol; por ello la marca no sustituye
la ramificación por una trayectoria aislada.

\Needspace{15\baselineskip}
Póngase \(R_m:=\mathbb Z/3^m\mathbb Z\), y sea
\(V_j(\zeta):=\operatorname{Desc}_{k,j}(\zeta)\). El transporte, los
proyectores de hoja y el residuo terminal se restringen al subárbol mediante
\[
 \mathcal U_j^\zeta:=
 \mathcal U_j\!\upharpoonright_{\ell^2(V_j(\zeta),\mu_j)},
 \qquad P_j^\zeta:=P_j\!\upharpoonright_{V_j(\zeta)},
 \qquad Q_j^\zeta:=I-P_j^\zeta,
\]
\[
 A_j^\zeta:=P_{j+1}^\zeta\mathcal U_j^\zeta P_j^\zeta
             +Q_{j+1}^\zeta\mathcal U_j^\zeta Q_j^\zeta,
 \qquad
 \eta_j^\zeta:=Q_{j+1}^\zeta\mathcal U_j^\zeta P_j^\zeta
             +P_{j+1}^\zeta\mathcal U_j^\zeta Q_j^\zeta.
\]
Para \(k\leq n\leq N\), sean
\[
 F_{k,k}^\zeta:=I,
 \qquad
 F_{k,n}^\zeta:=A_{n-1}^\zeta\cdots A_k^\zeta,
 \qquad
 T_{k,n}^\zeta:=(F_{k,n}^\zeta)^*F_{k,n}^\zeta.
\]
Los puertos propios del mismo subárbol son
\[
 \mathscr P_j^\Sigma(\zeta):=
 \{(y,\Sigma):y\in V_j(\zeta)\},
 \qquad
 \mathscr P_j^\Pi(\zeta):=
 \{(y,\Pi):y\in V_j(\zeta)\};
\]
denotamos por \(B_j^\zeta,\kappa_j^\zeta,\delta_j^\zeta\) las restricciones
de la incidencia y del complejo rectangular a esos puertos, y por
\(L_j^\zeta\) la precomposición inducida por la supresión de prefijos.

Para definir sin abreviaturas el lector de caminos, sea
\[
 W_\ell^\zeta(y):=
 R_m[\operatorname{Path}_{\mathscr T_{k,\ell}(\zeta)}
          (y,\operatorname{Term}_\ell)]
\]
y póngase
\[
 \mathscr D_{q,r;k,\ell}^{(m)}[\mathscr R]
 :=\bigoplus_{y_0\to\cdots\to y_q\ \text{en}\
                  \mathscr T_{k,\ell}(\zeta)}
 W_\ell^\zeta(y_q)\otimes_{R_m}
 \bigl(\Gamma_r(y_0)\otimes_{\mathbb Z}R_m\bigr),
 \qquad
 \operatorname{Tot}_d\mathscr D
 :=\bigoplus_{q+r=d}\mathscr D_{q,r;k,\ell}^{(m)}[\mathscr R].
\]
Las fórmulas ya demostradas producen, sobre cada
\(\mathscr R\in\mathfrak R_{k,N}(\zeta)\), los cinco objetos
\[
\begin{aligned}
 \mathcal F_{\rm sp}[\mathscr R]
 &:=(\ell^2(V_j(\zeta),\mu_j),
       \mathcal U_j^\zeta,A_j^\zeta,\eta_j^\zeta,
       T_{k,j}^\zeta)_{k\leq j\leq N},\\
 \mathcal F_{\rm sol}[\mathscr R]
 &:=(b,\kappa_{\rm can},C_\alpha,\kappa_\alpha)
      _{\alpha\subset\mathscr T_{k,N}(\zeta)},\\
 \mathcal F_{\rm gau}^{(m)}[\mathscr R]
 &:=(\mathscr P_j^\Sigma(\zeta),\mathscr P_j^\Pi(\zeta),
       B_j^\zeta,\kappa_j^\zeta,\delta_j^\zeta;R_m)_{k\leq j\leq N},\\
 \mathcal F_{\rm coh}^{(m)}[\mathscr R]
 &:=(\Gamma_y\otimes R_m,\partial,\Psi_\alpha,H_\bullet)
      _{y,\alpha\subset\mathscr T_{k,N}(\zeta)},\\
 \mathcal F_{\rm det}^{(m)}[\mathscr R]
 &:=(\operatorname{Det}(\operatorname{Tot}
       \mathscr D_{\bullet,\bullet;k,\ell}^{(m)}[\mathscr R]))
       _{k\leq\ell\leq N}.
\end{aligned}
\tag{C2}
\label{eq:cinco-lectores-comunes}
\]
Aquí \(\mathscr D_{\bullet,\bullet;k,\ell}^{(m)}[\mathscr R]\) es el
bicomplejo normalizado
de caminos y memoria de las secciones anteriores, restringido al subárbol de
horizonte \(\ell\) y reducido sobre \(R_m\). El último lector conserva la
colección ordenada de líneas determinantales hasta el horizonte \(N\); de este
modo, su truncamiento es una proyección canónica y no presupone una
factorización determinantal inexistente. La ley \(\mu\) es la ley canónica
del censo, y los puertos se fijan mediante la rama lexicográficamente mínima
del orden APP. Por tanto, ninguno de los cinco lectores recibe una ponderación
o un puerto externos.

Para
\(T\in\mathfrak T:=\{{\rm sp},{\rm sol},{\rm gau},{\rm coh},{\rm det}\}\),
defínase el espacio total tipado
\[
 \widetilde{\mathcal F}_{T;k,N}^{(m)}(\zeta)
 :=\coprod_{\mathscr R\in\mathfrak R_{k,N}(\zeta)}
   \{\mathcal F_T^{(m)}[\mathscr R]\}_{\mathscr R},
 \qquad
 p_T:\widetilde{\mathcal F}_{T;k,N}^{(m)}(\zeta)
      \longrightarrow\mathfrak R_{k,N}(\zeta),
\]
donde
\(\mathcal F_{\rm sp}^{(m)}:=\mathcal F_{\rm sp}\) y
\(\mathcal F_{\rm sol}^{(m)}:=\mathcal F_{\rm sol}\); es decir, el
superíndice \(m\) es inerte para esos dos lectores.
La proyección \(p_T\) olvida únicamente la estructura del lector y conserva
su registro fuente. El producto fibrado correcto es
\[
\begin{aligned}
 \mathcal J_{k,N}^{(m)}(\zeta):={}&
 \widetilde{\mathcal F}_{\rm sp;k,N}^{(m)}(\zeta)
 \mathop{\times}_{\mathfrak R_{k,N}(\zeta)}
 \widetilde{\mathcal F}_{\rm sol;k,N}^{(m)}(\zeta)
 \mathop{\times}_{\mathfrak R_{k,N}(\zeta)}
 \widetilde{\mathcal F}_{\rm gau;k,N}^{(m)}(\zeta)\\
 &\mathop{\times}_{\mathfrak R_{k,N}(\zeta)}
 \widetilde{\mathcal F}_{\rm coh;k,N}^{(m)}(\zeta)
 \mathop{\times}_{\mathfrak R_{k,N}(\zeta)}
 \widetilde{\mathcal F}_{\rm det;k,N}^{(m)}(\zeta).
\end{aligned}
\tag{C3}
\label{eq:constructor-correlativo}
\]
La igualdad de las cinco proyecciones obliga a conservar el mismo árbol, las
mismas aristas, la misma rama producida y el mismo registro. Para cada
\(\mathscr R\) las cinco construcciones determinan el elemento
\[
 s_{k,N}^{(m)}(\mathscr R)
 :=(\mathcal F_T^{(m)}[\mathscr R])_{T\in\mathfrak T}
 \in\mathcal J_{k,N}^{(m)}(\zeta).
\]
La existencia de este elemento procede de las operaciones de
\eqref{eq:cinco-lectores-comunes}, no de la
notación del producto fibrado.

La supresión de la última capa induce
\[
 \operatorname{res}^{\mathfrak R}_{N+1,N}
 \bigl(\mathscr R_{k,N+1}(\zeta;x)\bigr)
 :=\mathscr R_{k,N}(\zeta;\rho_N(x)).
\]
La restricción de cada operador a los generadores retenidos define
\[
 r^m_{N+1,N}:\mathcal J_{k,N+1}^{(m)}(\zeta)
 \longrightarrow\mathcal J_{k,N}^{(m)}(\zeta).
\]
En el lector determinantal, \(r^m_{N+1,N}\) suprime únicamente la última
coordenada de la colección de líneas de
\eqref{eq:cinco-lectores-comunes}. La reducción
\(R_{m+1}\twoheadrightarrow R_m\) define, por cambio de base,
\[
 q^N_{m+1,m}:\mathcal J_{k,N}^{(m+1)}(\zeta)
 \longrightarrow\mathcal J_{k,N}^{(m)}(\zeta).
\]
Sobre cada generador se cumplen
\[
\begin{gathered}
 r^m_{N+1,N}q^{N+1}_{m+1,m}
 =q^N_{m+1,m}r^{m+1}_{N+1,N},\\
 \partial\Psi_\alpha=\Psi_\alpha\partial,
 \qquad L_j^\zeta\kappa_j^\zeta
       =\kappa_{j+1}^\zeta L_j^\zeta,
 \qquad L_j^\zeta\delta_j^\zeta
       =\delta_{j+1}^\zeta L_j^\zeta,
\end{gathered}
\tag{C4}
\label{eq:naturalidad-D}
\]
y las composiciones de los mapas \(r\) y \(q\) satisfacen las identidades
functoriales. La primera igualdad se verifica porque borrar la última capa y
reducir los coeficientes actúan sobre coordenadas distintas. Las otras tres
son las identidades de naturalidad probadas anteriormente. Se obtiene un único
prosistema en \((N,m)\), no cinco límites yuxtapuestos.

En la raíz se define
\[
 \mathfrak J_{N,\bullet}^{(m)}:=\mathcal J_{0,N}^{(m)}(\ast).
\tag{C5}
\label{eq:fibras-conjuntas-horizonte}
\]
La estructura discreta conjunta del continuo es el límite de esos objetos
correlativos:
\[
 \boxed{
 \mathcal C_{\rm cont}^{\rm disc}
 :=\varprojlim_{\substack{N\geq1\\m\geq1}}
   \bigl(\mathfrak J_{N,\bullet}^{(m)},
         r^m_{N+1,N},q^N_{m+1,m}\bigr).}
\tag{C6}
\label{eq:continuo-discreto}
\]
Ésta es una definición por el prosistema generado por el árbol completo; no es
la imagen ni el grafo de una aplicación que retenga el estado de entrada como
primera coordenada.

Los registros poseen una marca terminal producida. Sus proyecciones compatibles
inducen
\[
 \Pi:\mathcal C_{\rm cont}^{\rm disc}\longrightarrow H_\infty^{\rm enr}.
\]
Recíprocamente, una historia completa determina, por sus prefijos ya emitidos,
las secciones compatibles \(s_{0,N}^{(m)}\). Así se obtiene, después de haber
definido el continuo,
\[
 \operatorname{Gen}_{\rm cont}:H_\infty^{\rm enr}
 \longrightarrow\mathcal C_{\rm cont}^{\rm disc},
 \qquad
 \Pi\circ\operatorname{Gen}_{\rm cont}
 =\operatorname{id}_{H_\infty^{\rm enr}}.
\tag{C7}
\label{eq:seccion-historias-continuo}
\]
Esta sección demuestra recuperabilidad; no se usa para definir
\(\mathcal C_{\rm cont}^{\rm disc}\).

\begin{theorem}[Emergencia correlativa de las cinco construcciones]
\label{thm:correlacion-cinco}
Para toda historia superviviente \(\zeta\in H_k\), todo \(N>k\), todo
\(m\geq1\) y todo \(\mathscr R\in\mathfrak R_{k,N}(\zeta)\), las cinco
construcciones de \eqref{eq:cinco-lectores-comunes} están definidas sobre el
mismo registro, y \eqref{eq:constructor-correlativo} contiene
\(s_{k,N}^{(m)}(\mathscr R)\). Sus truncamientos, reducciones modulares y
transportes satisfacen \eqref{eq:naturalidad-D}. Además,
\(\mathcal C_{\rm cont}^{\rm disc}\neq\varnothing\), la aplicación \(\Pi\)
es sobreyectiva y queda escindida por \(\operatorname{Gen}_{\rm cont}\). Las
cinco componentes son aspectos inseparables del mismo continuo, no cinco
dominios ni cinco aplicaciones independientes.
\end{theorem}

\begin{proof}
La supervivencia que define \(H_k\) y la ley de prolongación TPK dan
\(\rho_k^{-1}(\zeta)\neq\varnothing\) para cada vértice. Como las fibras son
finitas, el lema de Kőnig produce \(H_\infty^{\rm enr}\neq\varnothing\).

Fijado un registro \(\mathscr R\), el transporte de historias y hojas produce
\(\mathcal F_{\rm sp}[\mathscr R]\); la composición afín y el defecto exacto
de cancelación producen \(\mathcal F_{\rm sol}[\mathscr R]\); las incidencias
y la secuencia rectangular producen \(\mathcal F_{\rm gau}^{(m)}[\mathscr R]\);
el complejo de memoria y su transporte producen
\(\mathcal F_{\rm coh}^{(m)}[\mathscr R]\); y las líneas del ensamblaje
normalizado de caminos producen \(\mathcal F_{\rm det}^{(m)}[\mathscr R]\).
Cada construcción actúa sobre las mismas aristas generadas, no sobre un
subdominio supuesto ni sobre el ambiente abstracto de \(729\) palabras.

Las compatibilidades afín, de transporte, de incidencia y de memoria prueban
\eqref{eq:naturalidad-D}. Truncar la colección determinantal de
\eqref{eq:cinco-lectores-comunes} conmuta por
definición con \(r\); su cambio de base conmuta con \(q\). Las secciones
\(s_{0,N}^{(m)}\) son compatibles, de modo que producen
\(\operatorname{Gen}_{\rm cont}\). La identidad
\eqref{eq:seccion-historias-continuo} prueba a la vez la no
vacuidad del límite, la sobreyectividad de \(\Pi\) y la recuperabilidad de cada
historia sin introducirla como coordenada de definición.

La fibra, la relación, el número, la orientación, el acarreo, la memoria, la
frontera y la ruta permanecen en el registro común; la multisección, los
terminales y los lectores lineal y probabilístico permanecen en las cinco
construcciones correlativas. Por ello el paso al límite conserva la clausura de
dependencias y no disgrega el continuo.
\end{proof}

El reconocimiento ulterior del único continuo construido mediante lenguaje
espectral, solenoidal, de calibre, cohomológico o determinantal pertenece a la
comparación convencional posterior. No selecciona sus estados ni interviene en
su generación APP--TRIT--TPK.
